\section{Small degrees by exact character calculation}
\label{sec:finite}

The capacity calculations apply to the fixed large-degree candidate.
For the remaining small degrees we instead calculate exact character
residues, checking one self-conjugate candidate against every target.
A nonzero residue certifies a positive integer multiplicity; no estimate
or reconstruction of that integer is needed. We give the invariant of
the calculation and its implementation bounds. The source is reproduced
in Appendix~\ref{app:smallcode}.

\begin{proposition}\label{prop:finite-check}
For every integer $1\leq n\leq64$ other than $2,4,9$, there is a
self-conjugate partition $\lambda\vdash n$ such that
\[
  g(\lambda,\lambda,\nu)>0\qquad\text{for every }\nu\vdash n.
\]
\end{proposition}

\subsection{A recurrence for the entire coefficient vector}

Write $\chi^\lambda(\rho)$ for the character of $S^\lambda$ on the
conjugacy class of cycle type $\rho$. If $m_j(\rho)$ is the number of
parts of size $j$, put $z_\rho=\prod_j j^{m_j(\rho)}m_j(\rho)!$.
The character inner-product formula gives
\begin{equation}\label{eq:finite-character}
 g(\lambda,\lambda,\nu)
 =\sum_{\rho\vdash n}
   \frac{\chi^\lambda(\rho)^2\chi^\nu(\rho)}{z_\rho},
\end{equation}
since symmetric-group characters are real. The program computes this
sum for every $\nu\vdash n$ at once, recursively selecting the parts
of $\rho$ in decreasing order.

Calculations take place in $\mathbb F_p$, where
$p\in\{1000000007,1000000009\}$. Both numbers are prime, as is checked by
trial division through $31622$. Since $n\le64<p$, every denominator
in \eqref{eq:finite-character} is invertible modulo $p$.
All character values and vectors below are interpreted in this field.
For $s\geq0$, let $\mathcal P_s$ be the partitions of $s$, including
the empty partition when $s=0$. For a vector
$v=(v_\lambda)_{\lambda\in\mathcal P_s}$, write
\[
  \chi_v(\rho)=\sum_{\lambda\in\mathcal P_s}v_\lambda\chi^\lambda(\rho).
\]
Let $[\lambda]$ denote the coordinate vector with entry one at $\lambda$
and zero elsewhere, so $\chi_{[\lambda]}=\chi^\lambda$.

A border strip is an edge-connected skew diagram containing no $2\times2$
square; its leg length is one less than its number of occupied rows.  Let
$R_k$ denote the matrix that adds a border strip of size $k$, with entry
$(-1)^{\text{leg length}}$ for each permitted addition.  Its source and target
degrees will be understood from context. For $v$ of degree $s$ and
$\rho\vdash s-k$, the Murnaghan--Nakayama rule
\cite[Chapter~I, \S7, Example~5]{Macdonald}\cite[Rule~21.1]{James} gives
\begin{equation}\label{eq:finite-mn}
  \chi_{R_k^{\mathsf T}v}(\rho)
     =\chi_v((k)\cup\rho),
  \qquad
  \sum_\mu (R_k)_{\nu\mu}\chi^\mu(\rho)
     =\chi^\nu((k)\cup\rho).
\end{equation}

In a call \texttt{calc(s,k,e,v)}, the remaining cycle type has total
size $s$ and no part larger than $k$; $e$ counts the size-$k$ parts
already selected. The extra counter records the factorial in
$z_\rho$ when several equal parts are selected successively.
Reached states have $0\le s\le64$, $k\ge1$, and $e\ge0$.
Set $\rho_1=0$ for the empty partition and abbreviate $m_j=m_j(\rho)$.
Define $F_{s,k,e}(v)$ by its $\nu$ coordinate
\begin{equation}\label{eq:finite-invariant}
 F_{s,k,e}(v)_\nu
  =\sum_{\substack{\rho\vdash s\\\rho_1\leq k}}
   \chi_v(\rho)^2\chi^\nu(\rho)
   \frac{e!}{k^{m_k}(e+m_k)!}
   \prod_{1\leq j<k}\frac{1}{j^{m_j}m_j!}.
\end{equation}
Only states with $e+m_k\leq64$ are reached, so all displayed denominators
are invertible.  For $2\leq k\leq s$, separating the terms with $m_k=0$
from those with $m_k>0$ yields
\begin{equation}\label{eq:finite-recursion}
 F_{s,k,e}(v)
   =F_{s,k-1,0}(v)
     +\frac{1}{k(e+1)}R_k
        F_{s-k,k,e+1}(R_k^{\mathsf T}v).
\end{equation}
Indeed, \eqref{eq:finite-mn} supplies the three character factors, while for
$m=m_k>0$ the coefficient of the second term is
\[
 \frac{1}{k(e+1)}
 \frac{(e+1)!}{k^{m-1}(e+m)!}
   =\frac{e!}{k^m(e+m)!}.
\]
When $m_k=0$, the counter contributes only $e!/e!=1$, so the first
term resets it to zero. The counter is needed precisely when another
part of size $k$ is selected.
Each recursive call decreases $(s,k)$ in lexicographic order: either
the largest permitted part decreases, or the remaining degree decreases.
The recursion therefore terminates.

The base cases are exact.  At degree zero the answer is $v_\varnothing^2$.
If $0<s<k$, no size-$k$ part remains, so
$F_{s,k,e}(v)=F_{s,s,0}(v)$: the factor involving $e$ in
\eqref{eq:finite-invariant} is then $e!/e!=1$.  Calls with $k=1$ always have
$e=0$, since this case is returned immediately rather than recursed into.
Writing $d_s=(\dim S^\nu)_{\nu\vdash s}$ gives
\begin{equation}\label{eq:finite-base}
 F_{s,1,0}(v)=\frac{1}{s!}\,d_s(d_s^{\mathsf T}v)^2.
\end{equation}
The vector $d_s$ is computed by successive applications of $R_1$ to the
empty partition: the resulting Young-lattice chains are standard tableaux.
The array named \texttt{fac} in the source stores inverse factorials.
Finally, a zero vector $R_k^{\mathsf T}v$ contributes zero to the second
term of \eqref{eq:finite-recursion}; omitting that branch is therefore
valid even if the vector is zero only modulo $p$.

The recurrence and base cases prove by induction on $(s,k)$ that
\texttt{calc(s,k,e,v)} returns $F_{s,k,e}(v)$.
At the initial call $s=k=n$ and $e=0$, the weight in
\eqref{eq:finite-invariant} is $1/z_\rho$. Thus
\eqref{eq:finite-character} identifies the output as
\begin{equation}\label{eq:finite-output}
 F_{n,n,0}([\lambda])_\nu
   =g(\lambda,\lambda,\nu)\pmod p.
\end{equation}

\subsection{Exact implementation}

Pad a partition $\lambda$ with zero parts to length $64$, and encode it by
the beta set
\[
 B_\lambda=\{64+\lambda_j-j:1\leq j\leq64\}
          \subseteq\{0,\ldots,127\}.
\]
There are $64$ beads, and their positions sum to $2016+|\lambda|$.
The partition enumeration starts with $B_\varnothing=\{0,\ldots,63\}$
and moves the $j$th bead to $64+\lambda_j-j$ while selecting nonincreasing
parts.  It enumerates every partition exactly once.

Moving a bead from $x$ to an empty position $x+k$ adds a border strip of
size $k$.  If the bead passes $r$ other beads, the strip occupies $r+1$
rows, giving the sign $(-1)^r$.  For the bit mask $b$ of $B_\lambda$, the
expression
\[
 b\mathbin{\&}((\mathord{\sim}b)\mathbin{\gg}k)
\]
selects just the permitted starting positions.  Its top $k$ bits are zero,
so every selected position satisfies $x+k\leq127$.  With $m=2^x$ and
$t=2^{x+k}$, the mask $t-2m$ has exactly the bits strictly between the two
positions.  Thus \texttt{cnt((t-m-m)\&b)} supplies precisely the strip leg
length.  The routine \texttt{op} implements $R_k$ or its transpose according
to its direction argument.

The finite beta-set identity
\[
 B_{\lambda^t}=\{0,\ldots,127\}
                 \setminus\{127-x:x\in B_\lambda\}
\]
shows that \texttt{sym} tests self-conjugacy by requiring opposite bits at
positions $i$ and $127-i$.  Searching only these diagrams is legitimate;
indeed
\[
 g(\lambda,\lambda,(1^n))
   =\langle\chi^\lambda,\chi^\lambda\sgn\rangle
   =\begin{cases}1,&\lambda=\lambda^t,\\0,&\lambda\ne\lambda^t.\end{cases}
\]
The floating-point score orders candidates by the logarithm of their hook
product.  It is finite and affects only search order, never the arithmetic
certificate or the list of available candidates. The search stops as
soon as one candidate passes the test for every target.

We record the implementation bounds to make the use of machine arithmetic
explicit.  Residues lie in $[0,p-1]$; their sum is at most $2000000016$,
within a signed $32$-bit integer, and their product fits a signed $64$-bit
integer.  All mask shifts are between $0$ and $127$ and use unsigned
$128$-bit arithmetic.  The possible final enumeration argument
\texttt{pos=-1} occurs only after the last part has been placed, and the
routine returns before shifting.  Bit scans receive nonzero masks.
All degree indices lie in $[0,64]$.  In a non-base recursive step $k\geq2$,
at most $32$ parts of size $k$ can be selected, so the inverse-array indices
are also in range.

The partition-count recurrence, adding parts of sizes $1,2,\ldots,64$,
gives
\[
  |\mathcal P_{64}|=1741630,
  \qquad \sum_{s=0}^{64}|\mathcal P_s|=12308139.
\]
The largest padded hash table has $4194304$ entries and is at most half
full.  Hash arithmetic wraps only in unsigned types; collisions are
resolved by equality of complete $128$-bit masks.  To see that every queried
mask occurs in its table, sort its $64$ occupied positions as
$b_1>\cdots>b_{64}\geq0$.  The numbers $b_j-64+j$ are nonnegative and
nonincreasing, hence form a partition.  Moving one bead by $k$ increases
their sum by $k$, so the lookup is in precisely the enumerated target
degree.  Conversely, a size-at-most-$64$ partition never requires a
position outside the mask.

\begin{proof}[Proof of Proposition~\ref{prop:finite-check}]
For each self-conjugate candidate $\lambda$, the program evaluates
\eqref{eq:finite-output} modulo one or, if needed, both of the two
primes. Its array
\texttt{ok} is initialized afresh for that candidate, and records whether
each coordinate has a nonzero residue for at least one prime.  Therefore
a successful return for an eligible degree certifies one and the same
$\lambda$ for all target partitions $\nu$.  Since each Kronecker
coefficient is a nonnegative integer, a nonzero residue implies positive
multiplicity.  No bound on the coefficient sizes or reconstruction from
the two moduli is required.
A residue that vanishes for both primes is inconclusive: this is a
sufficient test for positivity, and an unsuccessful search would not
prove nonexistence.

The recorded calculation compiled the source in
Appendix~\ref{app:smallcode} with GNU C++~13.3.0 and executed it
for exactly the $61$ eligible degrees
$\{1,3,5,6,7,8\}\cup\{10,11,\ldots,64\}$.
Every invocation returned success. The complete execution record is
\path{verification/computation/character-61-degrees.log};
\path{verification/computation/RESULTS.json} records the source and log identities.
These successful invocations
establish the required existence statement in every eligible degree of
the finite range. The record contains exit statuses rather than the
selected partitions; their meaning is the successful return condition
of the source just analyzed.
\end{proof}

\subsection{Reproducibility and independent checks}

The files \texttt{smallcode\_submitted.cpp},
\texttt{smallcode\_dump.cpp}, and \texttt{smallcode\_verify.py} accompany
this proof under \path{verification/code/}; the execution records are under
\path{verification/computation/}. From \path{verification/code/}, the complete check can be
reproduced in Bash by
\begin{lstlisting}[language=bash,basicstyle=\ttfamily\footnotesize]
g++ -O3 -std=c++17 smallcode_submitted.cpp -o smallcode_submitted
python3 smallcode_verify.py runs 1 3 {5..8} {10..64} \
  --seconds 300 --log ../computation/recheck.log
\end{lstlisting}
Each child has a $1536$-MiB address-space limit, a $300$-second wall
limit and a $301$-second CPU limit. The driver returns failure if any
child fails or times out. The complete expected set of eligible degrees
must be checked: the three excluded degrees are deliberately skipped by
the C++ program and their zero exit statuses do not certify a square.

As an independent check of the implementation, the dump harness exposes
every square vector at degrees $1$ through $12$, together with dimensions,
self-conjugacy flags, and a nontrivial mixed input vector.  The Python
verifier computes integer characters by enumerating contained grid
diagrams and testing border strips for connectivity and the absence of a
$2\times2$ square.  Thus its strip enumeration and signs do not reuse the
bead implementation.  It checks character orthogonality, computes exact
integer Kronecker coefficients using conjugacy-class sizes, and compares
every exposed coefficient modulo both primes.  Every comparison passed,
including the dimensions and mixed vectors; the exact calculation also
finds no covering square in degrees $2,4,9$. The dump harness ran through
degree $12$ with undefined-behavior sanitization enabled and reported
no failure. These checks
are recorded in \path{verification/computation/character-independent-1-12.log}
and can be repeated with
\begin{lstlisting}[language=bash,basicstyle=\ttfamily\footnotesize]
g++ -O2 -std=c++17 -fsanitize=undefined \
  -fno-sanitize-recover=all smallcode_dump.cpp -o smallcode_dump_ubsan
python3 smallcode_verify.py compare {1..12} \
  --log ../computation/recheck-independent.log
\end{lstlisting}
Before comparing coefficients, the Python verifier requires every
expected record to occur exactly once, including every square vector
under both moduli. Its fault-injection tests reject missing or duplicate
records even when the child exits successfully. These tests can be run with
\begin{lstlisting}[language=bash,basicstyle=\ttfamily\footnotesize]
python3 -m unittest -v test_smallcode_verify.py
\end{lstlisting}
